\scanpage{26}
% Pencil foliation: 25, lower left.
\uuline{changement de base} $f^*$, \uuline{et compatible}
\uuline{également avec} $Rf_*$
[N.B. $T_\ell$ s'étend
évidemment en un foncteur
\diagram{\begin{tikzcd}
D^b(\mathcal{M}(X))\arrow[r,dash]&D^b(\mathcal{M}_\ell(X))
\end{tikzcd}}].
\unsure{$\{$} La détermination
des $T_\ell$ est encore ramenée au cas où
$X$ est de t. f. sur $\mathbb{Z}$.

\margin{compatible
avec $f^!$,
$Rf_!$, avec l'hom
\unsure{interne} etc \dots}

\margin{\uline{NB.} Ceci
exclut déjà
le choix \uline{trivial}
$\mathcal{M}^+(X)=0$ pour tt
$X$, car il
faudrait qu'on
ait $R^*f_*(\mathbb{Z}_\ell)=0$
pour $f,\ell$, ce
qui n'est
pas vrai
en général \dots}

\deleted{Si $X$ est un $\mathbb{Q}$-schéma
\inserted{de \deleted{\unsure{présentation}}\inserted{type} fini,}
alors on
aussi un foncteur
\diagram{\begin{tikzcd}
T_\infty=T_\infty^{(X)}:D^b(\mathcal{M}(X))\arrow[r]&\mathcal{M}_\infty(X)
\end{tikzcd}}
$\mathcal{M}_\infty(X)$ est la catégorie des \inserted{$\mathcal O_X$-}Modules
cohérents. Le foncteur $T_\infty$ est encore
\uline{exact} et \uline{fidèle}, compatible avec
$f^*$ et $Rf_*$.}

\margin{\uline{Signaler aussi la}
\uline{compatibilité de $T_\ell$ avec}
\uline{l'isom. de Künneth}}

5) Pour tout $X$, on a un \deleted{\unsure{un}}
élément canonique
\[
 \unsure{\mathbb{Q}|\mathbb{Z}}_X(-1)\ \text{ou }\mathbf1_X(-1)\in\operatorname{Ob}\mathcal{M}^+(X)
\]
